\begin{lemma}[A function and its classical gradient vanish off the topological support]
For $f\colon\mathrm{Plane}\to\mathbb{R}$ write $\operatorname{supp}f$ for the \emph{topological
support} of $f$, that is, the closure in $\mathrm{Plane}$ of the set
$\{x\in\mathrm{Plane}:f(x)\neq0\}$.  Then both of the following hold.
\begin{enumerate}
\item[(1)] For every $f\colon\mathrm{Plane}\to\mathbb{R}$ and every $x\in\mathrm{Plane}$, if
$x\notin\operatorname{supp}f$ then $f(x)=0$.
\item[(2)] For every $f\colon\mathrm{Plane}\to\mathbb{R}$ and every $x\in\mathrm{Plane}$, if
$x\notin\operatorname{supp}f$ then $Df(x)=0$, where $Df$ denotes the classical gradient field of
$f$.
\end{enumerate}

Each clause carries its own quantifiers over $f$ and over $x$, inside the conjunction, so the
statement has no hypotheses and no parameters of its own: a consumer obtains the pair once and
then applies each clause to whichever of its own test functions it needs.

Clause (2) is the one with content, and it is the reason the \emph{topological} support is the
right notion here.  Clause (1) says only that $f$ vanishes at a single point, which by itself
says nothing about derivatives; what clause (2) uses is that $\operatorname{supp}f$ is CLOSED,
so that its complement is open and $f$ vanishes on a whole neighbourhood of $x$.  Were
$\operatorname{supp}f$ replaced by the un-closed set $\{f\neq0\}$, the hypothesis
$x\notin\{f\neq0\}$ would again give $f(x)=0$, so clause (1) would survive, but clause (2) would
fail: $\{f=0\}$ need not be open, and a function can vanish at a point with nonzero gradient
there.
\end{lemma}
\begin{proof}
\emph{Clause (1).}  Let $f\colon\mathrm{Plane}\to\mathbb{R}$ and $x\in\mathrm{Plane}$ with
$x\notin\operatorname{supp}f$.  A set is contained in its own closure, so
$\{y:f(y)\neq0\}\subseteq\operatorname{supp}f$; contrapositively, since
$x\notin\operatorname{supp}f$ we get $x\notin\{y:f(y)\neq0\}$, that is, $f(x)=0$.

\medskip\noindent\emph{Clause (2).}  Let $f$ and $x$ be as given, with
$x\notin\operatorname{supp}f$, and put
\[
  V:=\mathrm{Plane}\setminus\operatorname{supp}f .
\]
The set $\operatorname{supp}f$ is a closure, hence closed, so $V$ is open; and $x\in V$ by
hypothesis.

Next, $f\equiv0$ on $V$: every $y\in V$ satisfies $y\notin\operatorname{supp}f$, so $f(y)=0$ by
clause (1), which has just been proved and whose statement is quantified over all points.

So $V$ is an open set on which $f$ vanishes identically.  Clause (3) of the background
proposition on the product rule and linearity for the classical gradient --- vanishing on an
open set --- therefore applies with this $V$ and gives that $f$ is differentiable at every point
of $V$ with $Df\equiv0$ on $V$.  Since $x\in V$, this yields $Df(x)=0$, which is clause (2).  As
$f$ and $x$ were arbitrary in each clause, both clauses hold as stated.
\end{proof}
